% Advisor report; provenance and reproduction are in README.md / SOURCES.md.
\input{preamble.tex}
\newcommand{\src}[1]{\note{Source: #1}}% sources moved to speaker notes (owner request 2026-10-14)
\newcommand{\epath}[1]{{\small\ttfamily\detokenize{#1}}}
\begin{document}
\raggedright

% 1
\begin{frame}[c]\relax
{\huge\bfseries\color{navy}WBR Perception\par}
{\Large\color{teal}Progress report · 13 October 2026\par}
\medskip{\color{teal}\hrule height .6pt}\medskip
\begin{columns}[c,onlytextwidth]
\begin{column}{.54\textwidth}
\Large Known baseline before\par camera adaptation\par\bigskip
\normalsize Pi 5 · D435i · ORB-SLAM3\par\medskip
\textbf{Accuracy evidence is encouraging.}\par
Physical qualification remains OPEN.
\end{column}
\begin{column}{.42\textwidth}
\slideimg[.43\textheight]{\linewidth}{P04}{D435i stereo camera}
\par\small\color{muted}Sensor illustration; no validation claim.
\end{column}\end{columns}
\note{About 15 minutes for slides 1--16; four backups follow. All results are historical evidence read on the Mac, not fresh experiments. No gate is promoted to PASS. The owner generates any new raster illustration; this report only reuses P04.}
\end{frame}

% 2
\begin{frame}{Goal: useful navigation without coupling balance to Linux}
\evidence{PLANNED}
\Large A wheel-biped robot needs scene-relative motion for navigation.\par\bigskip
\normalsize\deckitems{
\item Pi 5 estimates motion from stereo images and inertial data.
\item STM32H7 owns balance, hard-real-time control and motor safety.
\item Perception failure must remove optional navigation input safely.}
\vfill\takeaway{Qualify perception while preserving the controller ownership boundary.}
\src{AGENTS.md; docs/ARCHITECTURE.md}
\note{WBR means wheel-biped robot: torso, two articulated legs, one wheel at each foot. The boundary is a design requirement, not a tested robot-safety result. No mechanical drawing is claimed.}
\end{frame}

% 3
\begin{frame}{Ownership: the Pi provides an optional reference}
\evidence{PLANNED / OPEN integration}
\centering
\begin{tikzpicture}
\node[box,text width=71mm,minimum height=30mm,font=\normalsize,fill=navy!6] (pi) at (0,0) {\textbf{Raspberry Pi 5}\\Stereo + IMU acquisition\\Pose / map / health};
\node[box,text width=71mm,minimum height=30mm,font=\normalsize,fill=teal!8] (mcu) at (10,0) {\textbf{STM32H7}\\Body estimation + balance\\Watchdog + motor limits};
\draw[flow] (pi.east)--node[above,align=center,font=\small]{Optional\\reference}(mcu.west);
\node[box,text width=68mm,draw=fault,font=\normalsize] (fail) at (0,-3.2) {Pi stall / stale estimate};
\node[box,text width=68mm,draw=fault,font=\normalsize] (safe) at (10,-3.2) {Reject input; keep control local};
\draw[faultflow] (fail.east)--(safe.west);
\end{tikzpicture}
\vfill\takeaway{Receiver health semantics and physical fault tests are still future gates.}
\src{AGENTS.md; docs/ARCHITECTURE.md}
\note{This is the existing ownership model, not a new node graph or public message proposal. Transport, clock relation, age thresholds and fallback bounds require review in N6/N7. No motor or controller test was performed.}
\end{frame}

% 4
\begin{frame}{Data path: two ways to read the camera}
\evidence{Built (not yet run on the camera) / PLANNED}
\centering
\begin{tikzpicture}[font=\normalsize]
\node[box,text width=24mm,minimum height=12mm] (d) at (0,0) {D435i\\USB 3};
\node[box,text width=36mm,minimum height=12mm] (sdk) at (4.2,0) {librealsense\\\small Intel's camera library};
\node[box,text width=50mm,minimum height=12mm] (rec) at (9.8,1.0) {Built: own recorder\\\small writes a .db3 file};
\node[box,draw=muted,text width=50mm,minimum height=12mm] (ros) at (9.8,-1.0) {Planned: realsense-ros\\\small ROS 2 topics and bag (N3)};
\node[box,text width=40mm,minimum height=12mm] (orb) at (9.8,-3.1) {ORB-SLAM3 program\\\small offline file or live};
\node[box,text width=26mm,minimum height=12mm] (traj) at (4.2,-3.1) {Trajectory};
\node[box,text width=24mm,minimum height=12mm] (eval) at (0,-3.1) {Evaluation\\\small evo};
\draw[flow] (d)--(sdk);\draw[flow] (sdk.east)--(rec.west);\draw[flow,dashed] (sdk.east)--(ros.west);
\draw[flow] (rec.east)--++(0.5,0)|-(orb.east);
\draw[flow] (orb)--(traj);\draw[flow] (traj)--(eval);
\node[box,draw=muted,text width=30mm] (stm) at (9.8,-5.0) {Later: STM32H7};
\draw[flow,dashed] (orb)--node[right,font=\small]{health signal}(stm);
\end{tikzpicture}
\vfill\takeaway{Today's two programs use librealsense directly, without ROS 2; the ROS 2 path (N3) is still required.}
\src{Session facts §E; N4B prep/n4b\_plan.md}
\note{librealsense is Intel's RealSense SDK: the C++ library that talks to the camera over USB. The recorder and the D435i ORB-SLAM3 program built on 2026-10-13 use it directly (no ROS 2 node, no topics); the program reads either the recorded .db3 file (offline, N4B) or the camera (live, N4C). Both are built but have never run on the camera. The ROS 2 path (realsense-ros publishing topics, recorded with rosbag2) is the separate N3 gate and remains required for the robot integration.}
\end{frame}

% 5
\begin{frame}{Method: each gate needs artifacts and acceptance}
\evidence{PLANNED / OPEN}
\centering
\begin{tikzpicture}[font=\normalsize]
\foreach \x/\id/\txt in {0/a/N0 platform,4.5/b/N1 access,9/c/N2 raw,13.5/d/N3 ROS 2}{\node[box,text width=32mm,minimum height=12mm] (\id) at (\x,0) {\txt};}
\draw[flow](a)--(b);\draw[flow](b)--(c);\draw[flow](c)--(d);
\node[box,text width=35mm,minimum height=15mm] (ea) at (0,-2.5) {N4A EuRoC\\No camera};
\node[box,text width=35mm,minimum height=15mm] (eb) at (6.5,-2.5) {N4B D435i\\Offline};
\node[box,text width=35mm,minimum height=15mm] (ec) at (13.5,-2.5) {N4C D435i\\Live};
\draw[flow](a)--(ea);\draw[flow](d.south)--(13.5,-1.3)--(6.5,-1.3)--(eb.north);
\node[font=\small,anchor=south] at (10,-1.3) {Stereo + calibration};
\draw[flow](eb)--(ec);
\node[box,text width=125mm,minimum height=12mm] at (6.75,-4.6) {Qualified N4 trials $\rightarrow$ N5 metrics $\rightarrow$ N6 health $\rightarrow$ N7 boundary};
\end{tikzpicture}
\vfill\takeaway{Build success, plausible trajectories and historical receipts do not imply gate PASS.}
\note{N4A needs an eligible N0 and separate execution authorization; it does not depend on camera gates. N4B requires N3, a reviewed stereo extension and accepted calibration. N4C renews input qualification under stereo load. Gate definitions follow the course, sections 4.1--4.14; thresholds remain unchanged. N5--N7 await qualified inputs and reviewed bounds.}
\end{frame}

% 6
\begin{frame}{Gate status: the active unit needs fresh physical evidence}
\evidence{OPEN · as of 13 October 2026}
\begin{tabular}{@{}L{45mm}L{140mm}@{}}\toprule
\textbf{Gate} & \textbf{Current status}\\\midrule
N0 platform & OPEN: loaded power unstable\\[8pt]
N1 access & OPEN: new unit; prior clock stop\\[8pt]
N2 raw & OPEN: new-unit soak missing\\[8pt]
N3 ROS 2 & OPEN: record/replay unqualified\\\bottomrule
\end{tabular}
\vfill\takeaway{Old-unit N1/N2 and calibration receipts do not transfer to the new D435i.}
\note{Session facts supersede historical inventory drafts in PROJECT_STATE. New unit SDK 317422076650; USB descriptor 327743061057. Clock was corrected later but fresh sealed N1 evidence is still missing. N2 exact profile remains IR1 640x480 Y8/15, gyro 200, accel 100; mandatory 300-second soak unchanged.}
\end{frame}

% 7
\begin{frame}{Gate status: measured accuracy; adaptation still OPEN}
\evidence{OPEN / PLANNED · as of 13 October 2026}
\begin{tabular}{@{}L{45mm}L{140mm}@{}}\toprule
\textbf{Gate} & \textbf{Current status}\\\midrule
N4A EuRoC & OPEN: three runs; timing invalid\\[7pt]
N4B offline & OPEN: built; camera data absent\\[7pt]
N4C live & OPEN: built; never camera-run\\[7pt]
N5 metrics & PLANNED: full campaign pending\\[7pt]
N6 health & PLANNED: semantics/fixtures\\[7pt]
N7 boundary & PLANNED: bench review/tests\\\bottomrule
\end{tabular}
\note{Paths in the evidence table are provenance, not invitations to open reference decks during the meeting. N4A has preliminary ATE metrics, not qualified N5 campaign evidence. The latest build/run status comes from the authoritative session summary; earlier prep receipts describe Mac-only work and must not be mistaken for subsequent Pi runtime receipts.}
\end{frame}

% 8
\begin{frame}{Since the previous report: a runnable native baseline}
\evidence{MEASURED}
\begin{tabular}{@{}L{43mm}L{140mm}@{}}\toprule
\textbf{Changed} & \textbf{Consequence}\\\midrule
New D435i unit & Camera serial \texttt{317422076650}; USB serial \texttt{327743061057}.\newline Renew physical input and calibration evidence.\\[12pt]
Temporary Pi 5 & 4 GB RAM; 32 GB microSD; 5 V / 3 A supply.\newline Target 8 GB platform has not been measured here.\\[12pt]
Native ORB-SLAM3 & Upstream \texttt{4452a3c}; three compatibility patches.\newline Isolated Pi copy; EuRoC example and settings unchanged.\\\bottomrule
\end{tabular}
\vfill\takeaway{The algorithm path now runs; the hardware change restarts camera qualification.}
\src{Session facts §A--C; N4A/provenance.txt}
\note{New unit firmware 5.17.3.10, BMI085, USB 5000 Mbps. SDK, ASIC, firmware-update ID and descriptor serial are separate fields; USB/ASIC association here is reported evidence, not a general identity rule. Three patches: C++17/drop march native, disable failing Sophus tests, bool increment replaced with true. Pi build around 71 min at j2 with 4 GB temporary swap. No claim that the 8 GB target builds without swap; that remains inferred.}
\end{frame}

% 9
\begin{frame}{N4A setup: repeat a dataset with a known answer}
\evidence{MEASURED setup / SOURCE-SUPPORTED dataset}
\begin{columns}[T,onlytextwidth]\begin{column}{.46\textwidth}
\textbf{EuRoC MH01}\par\medskip
\deckitems{\item Public machine-hall drone sequence.
\item Stereo 20 Hz + IMU 200 Hz.
\item 184 s; 3,682 stereo frames.
\item Ground-truth body position.}
\end{column}\begin{column}{.50\textwidth}
\begin{tikzpicture}[node distance=8mm]
\node[box,text width=73mm,font=\normalsize] (a) {Pinned dataset + upstream YAML};
\node[box,text width=73mm,font=\normalsize,below=of a] (b) {Unmodified EuRoC example\\Three headless runs on Pi 5};
\node[box,text width=73mm,font=\normalsize,below=of b] (c) {Body trajectory vs ground truth\\evo · SE(3) alignment};
\draw[flow](a)--(b);\draw[flow](b)--(c);
\end{tikzpicture}
\end{column}\end{columns}
\vfill\takeaway{Known input separates native-build behavior from D435i acquisition problems.}
\src{N4A/REPORT.md; N4A/provenance.txt}
\note{EuRoC MH_01_easy from official ETH Machine Hall dataset archive. The SDK camera is not involved. Source library has compatibility patches; the stereo_inertial_euroc example, EuRoC.yaml and MH01 timestamps are unchanged. Compare body/IMU exports with body ground truth, after ns-to-s conversion. Association tolerance 0.02 s. Ground truth is laser tracker/motion capture dataset evidence, not a reference collected by this project.}
\end{frame}

% 10
\begin{frame}{N4A result: 0.037--0.042 m translation ATE}
\evidence{MEASURED runs / SOURCE-SUPPORTED paper comparison}
\begin{columns}[T,onlytextwidth]\begin{column}{.38\textwidth}
\begin{tabular}{@{}lr@{}}\toprule
\textbf{Trial} & \textbf{RMSE (m)}\\\midrule
Run 1 & 0.0415\\[6pt]
Run 2 & 0.0368\\[6pt]
Run 3 & 0.0416\\[6pt]
\midrule Mean & 0.040\\[6pt]
Paper MH01 & 0.036\\\bottomrule
\end{tabular}\par\medskip
\small Paper: median of 10 runs.\par
Our three-run median: 0.0415 m.
\end{column}\begin{column}{.58\textwidth}
\begin{tikzpicture}[x=1cm,y=70cm,font=\small]
\foreach \v/\l in {0/0,0.01/0.01,0.02/0.02,0.03/0.03,0.04/0.04,0.05/0.05}{\draw[rule](0,\v)--(8.2,\v);\node[anchor=east]at(-.15,\v){\l};}
\draw[navy](0,0)--(8.2,0);\draw[navy](0,0)--(0,.052);
\foreach \x/\v/\l in {1.5/0.0415/Run 1,4/0.0368/Run 2,6.5/0.0416/Run 3}{\fill[teal](\x-.55,0)rectangle(\x+.55,\v);\node[above]at(\x,\v){\v};\node[below]at(\x,0){\l};}
\draw[fault,dashed,thick](0,.036)--(8.2,.036);
\draw[fault,dashed,thick](.1,-.014)--(.9,-.014);
\node[anchor=west,text=fault]at(1,-.014){Paper: 0.036 m};
\node[anchor=west]at(0,.058){ATE RMSE (m); lower is better};
\end{tikzpicture}
\end{column}\end{columns}
\vfill\takeaway{All three runs: IMU initialized, one map, no resets; N4A remains OPEN.}
\src{N4A/eval/ape\_summary.txt; \href{https://arxiv.org/pdf/2007.11898}{Campos et al., 2021, Table II}}
\note{Paper value verified in accepted author manuscript arXiv v2, PDF page 12, Table II, stereo-inertial row, MH01 column. Paper reports median of ten executions. This is similar accuracy on one sequence, not exact statistical reproduction. Local raw RMSE 0.041461, 0.036844, 0.041604 m. Rounded mean 0.040 m. No Sim(3) scale correction is used in these results. Bars start at zero. Under-voltage excludes valid performance timing.}
\end{frame}

% 11
\begin{frame}{N4A supports the native build, with a narrow scope}
\evidence{INFERRED from MEASURED runs / OPEN}
\begin{columns}[T,onlytextwidth]\begin{column}{.48\textwidth}
\begin{block}{Supported on MH01}
\deckitems{\item ARM64 build produces repeat trajectories.
\item Compatibility patches retain similar accuracy.
\item Initialization and map continuity are observed.}
\end{block}
\end{column}\begin{column}{.48\textwidth}
\begin{block}{Still unproven}
\deckitems{\item New D435i pixels and motion delivery.
\item Camera--IMU timing and calibration.
\item Real-time speed or robot navigation.}
\end{block}
\end{column}\end{columns}
\vfill\takeaway{Retain the known-data baseline while qualifying each camera-specific input.}
\src{N4A/REPORT.md; Session facts §D}
\note{One sequence and three runs do not prove broad generalization or exact equivalence with the paper. No RPE or wider sequence campaign was reported. Peak RSS 823--828 MiB; this is SLAM runtime memory, distinct from the native compile memory pressure. Wall time includes vocabulary loading and was recorded under power throttling.}
\end{frame}

% 12
\begin{frame}{D435i pipeline: built and linked; never camera-run}
\evidence{MEASURED build, session-reported / OPEN execution}
\begin{tikzpicture}[font=\normalsize]
\node[box,text width=48mm,minimum height=20mm] (r) at (0,0) {\texttt{d435i\_record}\\Stereo + raw motion};
\node[box,text width=38mm,minimum height=20mm] (bag) at (6,0) {librealsense recording\\+ timestamp CSV};
\node[box,text width=56mm,minimum height=20mm] (driver) at (12.5,0) {Patched stereo-inertial\\D435i example};
\draw[flow](r)--(bag);\draw[flow](bag)--node[above,font=\small]{\texttt{--bag}}(driver);
\node[below=10mm of driver,align=center,font=\small,text=navy]{No \texttt{--bag}: later live mode\\\texttt{--no-viewer}: headless};
\end{tikzpicture}\par\bigskip
IR1 + IR2: 640$\times$480 Y8 / 30 Hz; gyro 200 Hz; accel 100 Hz.\par\medskip
Emitter OFF. Accel interpolated onto gyro times; invalid time steps rejected.
\vfill\takeaway{Factory geometry is nominal; IMU noise is ASSUMED; calibration remains OPEN.}
\src{Session facts §E; N4B prep source/configuration}
\note{Session summary reports both programs built in ~/wbr-runtime/n4b-20261013T080913Z/bin with clean linkage. Local Mac verification.md predates that Pi build and is tooling-only; no fresh Pi build receipt was collected for this deck. Recorder writes stereo_imu.db3 plus timestamps.csv; confirm actual installed-SDK container/playback semantics on first gated take. This N4 30 Hz stereo profile is separate from N2's IR1-only 15 Hz baseline. Retain SDK warning: IMU Calibration is not available. Configuration fx=fy 386.95, cx 321.03, cy 241.04, baseline 0.0499 m comes from nominal factory evidence; it is not a calibrated VIO claim.}
\end{frame}

% 13
\begin{frame}{Blockers: power first, then trustworthy input}
\evidence{Updated 14 October 2026 · 8 GB Pi now in use}
{\small\begin{tabular}{@{}L{28mm}L{80mm}L{76mm}@{}}\toprule
\textbf{Risk} & \textbf{Evidence and consequence} & \textbf{What can be done now}\\\midrule
Power / timing & 4 GB Pi: under-voltage in 49\% of samples; timing invalid. 8 GB Pi: under-voltage within 3 min of boot. & \textbf{Not fixed.} Interim: shorter, thicker USB-C cable or powered hub for the D435i. Real fix: 5 A supply.\\[5pt]
Storage & Old card: worst write 17 MB/s vs 18.4 MB/s stereo stream; no margin. & \textbf{Waiting} for the 128 GB card: clone, then a sustained-write test of at least 30 MB/s.\\[5pt]
Calibration & Nominal camera--IMU geometry; IMU noise values copied, not measured. & \textbf{Partly.} Prepare data collection now; the 4-hour IMU noise recording needs clean power.\\[5pt]
RAM & 4 GB Pi needed 4 GB swap to build ORB-SLAM3. & \textbf{Fixed by the 8 GB Pi} (7.8 GiB, no swap). EuRoC rerun done; numbers pending.\\\bottomrule
\end{tabular}}
\src{Session facts §A,C--F; N4A/REPORT.md and power\_window.log; Pi setup thread findings 2026-10-14}
\note{8 GB Pi findings were relayed from the Pi setup and ORB-SLAM3 thread on 2026-10-14: 7.8 GiB RAM, no swap; MH01 run completed, its RAM and ATE numbers not yet reported; under-voltage within 3 minutes of boot. Ordered 5.1 V/5 A supply and 128 GB card are plans, not measured fixes.}
\end{frame}

% 14
\begin{frame}{Next steps: restore prerequisites in order}
\evidence{PLANNED / OPEN exit criteria}
\begin{tabular}{@{}L{9mm}L{55mm}L{121mm}@{}}\toprule
 & \textbf{Action} & \textbf{Exit evidence}\\\midrule
1 & Stabilize platform / storage & Fresh N0: synced UTC, flags 0x0; qualified recording storage.\newline Repeat known-data timing under healthy loaded power.\\[12pt]
2 & Renew new-unit N1 & Camera serial + own USB branch $\geq$5000M + normal-user access.\newline Fresh receipt after storage/topology changes.\\[12pt]
3 & Complete N2 $\rightarrow$ N3 & N2: 20 $\rightarrow$ 105 $\rightarrow$ mandatory 300 s; unchanged thresholds.\newline N3: stream/rate/time/QoS evidence, finalized bag, isolated replay.\\\bottomrule
\end{tabular}
\vfill\takeaway{Each step starts only when the previous step's exit evidence exists.}
\src{Phase A runbook; N4B prep/n4b\_plan.md}
\note{No physical command is executed by this slide build. N2 exact profile IR1 640x480 Y8/15, gyro/200, accel/100. N2 thresholds IR 14.25--15.75, gyro 190--210, accel 95--105 Hz; image gap <=0.20 s, motion gap <=0.05 s, positive stamps/frame steps, no missing image frames, clean logs/content/human review. N0 temperature <75 C and available memory >=512 MiB. Large datasets need separately qualified storage; a larger microSD alone is not qualification. Keep acquisition local to Pi; SSH control only.}
\end{frame}

% 15
\begin{frame}{Then: immutable offline trials before live timing}
\evidence{PLANNED / OPEN exit criteria}
\begin{tabular}{@{}L{9mm}L{55mm}L{121mm}@{}}\toprule
 & \textbf{Action} & \textbf{Exit evidence}\\\midrule
4 & Review stereo + calibration & Attest paired IR1/IR2 at 30 Hz, IMU coverage and storage load.\newline Qualify geometry, noise and time offset; preserve raw warnings.\\[12pt]
5 & N4B offline D435i & Hash the take/settings; reconcile pixels, stamps and counts.\newline Repeat initialization, finite trajectory and map-continuity checks.\\[12pt]
6 & N4C live $\rightarrow$ N5 & Renew stereo input; measure age, queue growth, losses and health.\newline Review performance bounds; accuracy requires valid reference.\\\bottomrule
\end{tabular}
\vfill\takeaway{Nominal-config development runs can inform diagnosis; they cannot qualify VIO.}
\src{N4B prep/n4b\_plan.md; docs/ORB\_SLAM3\_OFFLINE\_PREPARATION.md}
\note{The first planned guided development take is 60--90 s with translation and rotations, a loop and a still ending. It is conditional on preceding gates and separately qualified storage. Initialization should appear in the log; inspect full atlas/state history, as largest-map exports can hide losses. Development criteria are not approved numeric acceptance budgets. Review N6 health semantics before separately authorized N7 boundary tests. No unsupported request to the advisor or invented procurement requirement is added.}
\end{frame}

% 16
\begin{frame}{Summary: keep the baseline; qualify the camera next}
\evidence{MEASURED progress / OPEN qualification / PLANNED next work}
\Large Three EuRoC runs: 0.037--0.042 m ATE.\par\bigskip
\normalsize\deckitems{
\item Native Pi build is usable on the known-data sequence.
\item D435i recorder and adapter are built, but never camera-run.
\item Resolve power/storage, then renew N1 $\rightarrow$ N2 $\rightarrow$ N3.
\item Offline stereo/calibration qualification precedes live performance.}
\vfill\takeaway{No gate PASS is claimed. STM32H7 retains balance and motor safety.}
\par\medskip{\small\color{muted}Further detail: course reference deck and learning deck.}
\note{Close the meeting here; four backup slides support questions about alignment, patches, power and storage. Advisor decisions on metric budgets/calibration remain future reviews; no deadline or unsupported request is invented.}
\end{frame}

\appendix
% 17
\begin{frame}{Backup: ATE uses one rigid alignment, without scale}
\evidence{MEASURED evaluation method}
Absolute trajectory error: matched body-position differences after alignment.\par
\[
(R^*,t^*)=\arg\min_{R\in SO(3),\,t}\sum_{i=1}^{n}\|Rp_i+t-q_i\|^2,
\qquad
\mathrm{ATE}_{\mathrm{RMSE}}=\sqrt{\frac{1}{n}\sum_{i=1}^{n}\|R^*p_i+t^*-q_i\|^2}
\]
{\small $p_i$: estimate; $q_i$: ground truth; $n$: matched positions.}\par\medskip
\centering\begin{tikzpicture}
\node[box,text width=47mm,font=\normalsize] (a) at (0,0) {Body poses\\ns $\rightarrow$ s};
\node[box,text width=48mm,font=\normalsize] (b) at (6.2,0) {Associate to GT\\Tolerance 0.02 s};
\node[box,text width=48mm,font=\normalsize] (c) at (12.4,0) {Align SE(3)\\Position RMSE};
\draw[flow](a)--(b);\draw[flow](b)--(c);
\end{tikzpicture}\par\bigskip
\raggedright\texttt{evo\_ape tum gt.tum est.tum -a --t\_max\_diff 0.02}
\vfill\takeaway{Stereo-inertial has metric scale; Sim(3) would answer a different question.}
\src{N4A/REPORT.md; N4A/eval/ape\_summary.txt}
\note{p_i estimate, q_i associated ground truth; both body positions. R and t are one global alignment, not local per-segment fits. n is associated pose count; retained exports have 3678 poses, but associated count is not re-invented. Translation ATE is in metres. RPE, alignment residuals, coverage exclusions and full metric campaign remain incomplete.}
\end{frame}

% 18
\begin{frame}{Backup: three compatibility patches in an isolated copy}
\evidence{MEASURED build / INFERRED narrow behavioral consistency}
\begin{tabular}{@{}L{53mm}L{128mm}@{}}\toprule
\textbf{Patch} & \textbf{Reason / limit}\\\midrule
C++17; remove\newline\texttt{-march=native} & Pangolin 0.9.1 / Ubuntu 24.04 toolchain compatibility.\newline Compile configuration changed; retain source and binary hashes.\\[13pt]
Sophus tests OFF & GCC 13 aarch64 array-bounds warning fails their \texttt{-Werror}.\newline ORB-SLAM3 uses Sophus headers; disabled tests are not validated.\\[13pt]
\texttt{mnFullBAIdx++}\newline$\rightarrow$ \texttt{= true} & C++17 forbids incrementing bool.\newline Preserves the prior C++11 bool-increment effect.\\\bottomrule
\end{tabular}
\vfill\takeaway{MH01 accuracy is evidence for this build; it does not validate every code path.}
\src{Session facts §C; Pi build/build\_orbslam3.sh; N4A/provenance.txt}
\note{Upstream 4452a3c4ab75b1cde34e5505a36ec3f9edcdc4c4. Build copied into ~/orbslam3/ORB_SLAM3, repo checkout preserved. Full build logs/patches are partly Pi-side, as documented in the session facts; do not imply all raw logs are on the Mac. Example and EuRoC configuration remained upstream. No fresh build or source mutation occurred for this deck.}
\end{frame}

% 19
\begin{frame}{Backup: live power faults in 87 of 178 samples}
\evidence{MEASURED samples / CALCULATED 48.9\%}
\centering\begin{tikzpicture}[font=\normalsize,x=.085cm]
\fill[fault](0,0)rectangle(87,1);
\fill[teal!45](87,0)rectangle(178,1);
\node[above,align=center]at(43.5,1.2){87 samples\\\texttt{0x50005}};
\node[above,align=center]at(132.5,1.2){91 samples\\\texttt{0x50000}};
\node[below,align=center,text=fault]at(43.5,-.2){Live faults + history};
\node[below,align=center]at(132.5,-.2){Historical flags remain};
\end{tikzpicture}\par\bigskip
\raggedright 5-second sampling during the EuRoC runs.\par\medskip
Maximum observed temperature: about 63\,$^{\circ}$C.\par\medskip
Pi UTC labels were about 80 minutes behind; NTP corrected later.
\vfill\takeaway{Historical-only flags are not a clean N0 result; timing must be repeated.}
\src{N4A/power\_window.log; N4A/REPORT.md clock note}
\note{0x50005 records live under-voltage and throttling plus historical bits; 0x50000 retains the historical bits. No claim about exact duty cycle between samples. Maximum logged temperature 63.1 C. Build had 46 events per session summary, separate from this run sample denominator. Clock offset affects cross-machine UTC correlation; within-run elapsed durations have no observed clock step, but wall time includes vocabulary load and is not per-frame compute.}
\end{frame}

% 20
\begin{frame}{Backup: stereo payload already consumes the old write margin}
\evidence{CALCULATED payload / MEASURED historical probe / PLANNED capacity}
\[
640\times480\times1\times2\times30
=18{,}432{,}000\ \mathrm{B/s}
\approx18.4\ \mathrm{MB/s}
\]
\begin{tabular}{@{}L{62mm}L{119mm}@{}}\toprule
\textbf{Quantity} & \textbf{Implication}\\\midrule
60--90 s stereo payload & About 1.11--1.66 GB per take, before recording overhead.\\[10pt]
Old write probe: 17--19 MB/s & Worst case below payload; nominal result has little margin.\\[10pt]
32 GB card: 13 GB free & Capacity alone does not prove sustained recording throughput.\\[10pt]
128 GB: planned 95 GB data & Roughly 55--90 takes in session estimate; not qualified storage.\\\bottomrule
\end{tabular}
\vfill\takeaway{Qualify storage under load and copy takes to the Mac after each session.}
\src{Session facts §F; N4B prep storage budget}
\note{Decimal units: MB=10^6 B, GB=10^9 B. Calculations include only two Y8 image planes, not IMU, metadata, SQLite/container overhead, filesystem or headroom. Exact image-only 95 GB / 1.65888 to 1.10592 GB gives about 57--86 full takes; report retains the session's rounded estimate 55--90 and labels it planned. V30/A2 card ordered, not received/qualified. MicroSD policy remains OS/tools/small logs; large captures require separately qualified storage. Preserve historical logs; proposed card clone and migration are not executed here.}
\end{frame}
\begin{frame}{References}
\small
C. Campos, R. Elvira, J. J. G\'omez Rodr\'iguez, J. M. M. Montiel, J. D. Tard\'os. ORB-SLAM3: An accurate open-source library for visual, visual--inertial and multi-map SLAM. \emph{IEEE Transactions on Robotics}, 37(6), 2021. \href{https://arxiv.org/abs/2007.11898}{arXiv:2007.11898}\par\medskip
M. Burri et al. The EuRoC micro aerial vehicle datasets. \emph{International Journal of Robotics Research}, 35(10), 2016.\par\medskip
Project evidence (experiment reports, logs and trajectories) is listed per slide in the speaker notes.
\end{frame}
\end{document}
